\documentclass[11pt]{article}

\usepackage[a4paper,margin=27mm]{geometry}
\usepackage{amsmath,amssymb,amsthm,mathtools}
\usepackage{enumitem}
\usepackage{hyperref}

\hypersetup{
  colorlinks=true,
  linkcolor=blue!50!black,
  citecolor=blue!50!black,
  urlcolor=blue!50!black
}

\setlength{\emergencystretch}{2em}

\theoremstyle{definition}
\newtheorem{definition}{Definition}
\newtheorem{remark}[definition]{Remark}
\newtheorem{convention}[definition]{Convention}

\theoremstyle{plain}
\newtheorem{proposition}[definition]{Proposition}

\DeclareMathOperator{\supp}{supp}
\DeclareMathOperator{\Read}{Read}
\DeclareMathOperator{\Reg}{Reg}
\newcommand{\qu}{\mathrm u}
\newcommand{\qd}{\mathrm d}
\newcommand{\qc}{\mathrm c}
\newcommand{\qs}{\mathrm s}
\newcommand{\qt}{\mathrm t}
\newcommand{\qb}{\mathrm b}

\title{The Grand Meson}
\author{Shuhei Ihara}
\date{\today}

\begin{document}
\maketitle

\begin{abstract}
This paper introduces the Grand Meson, not as a theory of everything, but
as a safe \(n=3\) general-purpose mesonic object.  The starting point is the
General Meson: a two-port Yukawa-style mediator read, in modern quark
language, through the light pair \(\qu\leftrightarrow\qd\).  The Grand
Meson opens the two remaining readable quark-pair axes and thereby becomes
a three-axis, six-port mediator.  The extension is conservative: the
original two-port meson is recovered by projection, while the new ports are
registered as defined, null, or active instead of being left outside the
type of the theory.  The central claim is modest and structural: the Grand
Meson is the first complete and last safe mesonic case for a
three-dimensional observer.
\end{abstract}

\section{Opening}

A meson is usually approached through its physical role: it mediates an
interaction.  In the present paper the same word is first read through its
paired role.  A meson stands between two complementary light-flavor faces.
In the Yukawa-facing reconstruction used here, that pair is
\[
  \qu\leftrightarrow\qd.
\]
This pair is not yet a full dynamics.  It is the minimum shape of a
mediator: an oriented reading may run \(\qu\to\qd\) or \(\qd\to\qu\), and
the middle term makes the relation nameable.  Yukawa did not formulate his
meson with quark flavors; the notation is a present reconstruction of the
light pair behind the proton--neutron mediation problem.

The former working nickname ``Ethon'' is retired here.  The technical
object of this paper is the Grand Meson.  The name is structural, not
honorific: ``grand'' means that a two-port light-pair meson has been closed
into the safe three-axis case.

\section{The General Meson}

Let the Yukawa-facing readable pair be
\[
  A_Y=\{\qu,\qd\}.
\]
Let \(\mathsf{Pol}\) be a class of admissible mesonic polymers.  The
smallest implementation may be generated from
\[
  \mathbb P_3=\{\mathbf{-1},\mathbf 0,\mathbf{+1}\},
\]
but no physical charge, amplitude, or particle identity is assigned to
these pseudo-values at this level.

\begin{definition}[General Meson]
A General Meson is a two-port mesonic record
\[
  m:A_Y\to\mathsf{Pol}.
\]
Equivalently,
\[
  \mathsf{Gen}:=\mathsf{Pol}^{A_Y}.
\]
The ports \(\qu\) and \(\qd\) are the two light-flavor faces of the
Yukawa-style mesonic record.
\end{definition}

This definition deliberately keeps the Yukawa-style mediator idea narrow.
It does not say that a meson is universal.  It says only that the basic
mesonic act is readable mediation across the \(\qu,\qd\) light pair.  In
that sense a General Meson is the optimal eye for the nuclear-force
problem: the proton--neutron transition axis is already specified by the
problem itself.

\section{The Hidden Cost of Undefined Ports}

The decisive distinction is between null and undefined.  A null port is
known to be empty.  An undefined port has not entered the type of the
theory.

\[
  \mathbf 0=\text{registered null},
  \qquad
  \bot=\text{unregistered undefined}.
\]

\begin{definition}[Undefined lift]
Let
\[
  F_6=\{\qu,\qd,\qc,\qs,\qt,\qb\},
  \qquad
  \mathsf{Pol}_{\bot}:=\mathsf{Pol}\cup\{\bot\}.
\]
The undefined lift of a General Meson \(m\in\mathsf{Gen}\) is
\[
  L_{\bot}(m):F_6\to\mathsf{Pol}_{\bot}
\]
defined by
\[
  L_{\bot}(m)(\qu)=m(\qu),\qquad L_{\bot}(m)(\qd)=m(\qd),
\]
and
\[
  L_{\bot}(m)(q)=\bot
  \quad(q\in F_6\setminus A_Y).
\]
\end{definition}

Thus the remaining axes are not false, not zero, and not disproved.  They
are absent from the original type.  The two-port theory remains locally
correct while it cannot register consequences that travel through the
untyped directions.

\begin{remark}[Meta-theoretic reading]
The undefined ports are not a defect in the original General Meson.  They
are the price of a successful restriction.  A narrow mediator may solve its
local problem precisely because it does not carry the full burden of every
other axis.
\end{remark}

\section{The Grand Meson}

The Grand Meson is obtained by replacing undefined ports with readable
ports.  A port may be active, or it may be null, but it may not be outside
the type.

\begin{definition}[Grand Meson]
A Grand Meson is a six-port mesonic record
\[
  g:F_6\to\mathsf{Pol}.
\]
Equivalently,
\[
  \mathsf{Grand}:=\mathsf{Pol}^{F_6}.
\]
The six ports are grouped into three readable axes:
\[
  \qu\leftrightarrow\qd,\qquad
  \qc\leftrightarrow\qs,\qquad
  \qt\leftrightarrow\qb.
\]
\end{definition}

In this notation the Grand Meson is the safe \(n=3\) object:
\[
  \mathsf{Grand}=\mathsf{Pol}^{2\cdot 3}.
\]
The adjective ``safe'' means that all three axes are typed.  It does not
mean that every transition is allowed, nor that every physical theory has
been solved.

\begin{definition}[Registered support]
For \(g\in\mathsf{Grand}\), define
\[
  \supp(g)=\{q\in F_6\mid g(q)\ne\mathbf 0\}.
\]
A Grand Meson may have inactive ports, but inactivity is registered by
\(\mathbf 0\), not by \(\bot\).
\end{definition}

\section{Conservative Recovery}

The Grand Meson must not destroy the General Meson.  It must contain it as
the two-port face that was already working.

\begin{definition}[Projection to the General Meson]
Define the \(\qu,\qd\)-projection
\[
  p_{\qu\qd}:\mathsf{Grand}\to\mathsf{Gen}
\]
by
\[
  p_{\qu\qd}(g)=g|_{\{\qu,\qd\}}.
\]
\end{definition}

\begin{proposition}[Conservative recovery]
Every General Meson is recovered from any Grand Meson whose \(\qu,\qd\)-face
agrees with it.  In particular, if \(g(\qu)=m(\qu)\) and \(g(\qd)=m(\qd)\), then
\[
  p_{\qu\qd}(g)=m.
\]
\end{proposition}

\begin{proof}
The projection \(p_{\qu\qd}\) discards the four ports outside
\(\{\qu,\qd\}\) and keeps exactly the light-pair values.  Therefore
agreement on those two ports is sufficient and necessary for recovery.
\end{proof}

This is the main discipline of the construction.  The Grand Meson does not
rewrite the older meson.  It keeps the working \(\qu\leftrightarrow\qd\)
face and opens two additional quark-pair faces under explicit registration.

\section{Einstein's Gate}

The Grand Meson must also avoid the classical aether failure.  It is not a
hidden medium with an absolute state of rest.  It is a local six-port
mediator whose outputs must be readable without introducing an absolute
frame.
Einstein's Gate is not a claim of Einstein's endorsement.  It names the
admissibility condition inherited from relativity: no mesonic theory may
reintroduce an absolute hidden medium.

\begin{convention}[No absolute mesonic frame]
No equation in the Grand Meson theory may depend on a privileged global
rest frame of the mediator.  Only typed ports, local registration, and
observer-readable projections may be used.
\end{convention}

This convention is Einstein's Gate.  It does not ban mediation.  It bans
absolute mediation.  The Grand Meson is therefore not a return to the
classical aether.  It is a typed local mediator that must survive
projection into admissible observations.

\section{Why \texorpdfstring{\(n=3\)}{n=3} Is the Safe Case}

Formally one may define an \(n\)-meson by
\[
  \mathsf{Mes}_n:=\mathsf{Pol}^{2n}.
\]
The case \(n=1\) is the General Meson.  The case \(n=3\) is the Grand
Meson.

The reason to stop at \(n=3\) is not that higher \(n\) is impossible.  It
is that a three-dimensional observer already has the first complete
rotational reading at \(n=3\).  Three axes close the ordinary posture
problem:
\[
  \Read_{O_3}(\mathsf{Mes}_3)
  =
  \Read_{O_3}(\mathsf{Pol}^{6}).
\]
For \(n>3\), additional axes require additional readout structure.  Such
structure may be meaningful in a later theory, but it is not required for
the present observer class:
\[
  \Read_{O_3}(\mathsf{Mes}_n)
  \not>
  \Read_{O_3}(\mathsf{Mes}_3)
  \qquad(n>3).
\]

Thus \(n=3\) is the first complete and last safe case: complete because
the three axes are readable as a closed posture, safe because the theory
does not ask the reader to chase an infinite ascent of unread axes.

\section{Safe Transition}

Safety is not the absence of change.  Safety is the registration of change.
Let a transition be written
\[
  u:g\to g'.
\]
The transition may change ports, support, or observable posture.  It is
unsafe only when the change crosses a typed boundary without registration.

\begin{definition}[Grand-safe transition]
Let \(\Reg(u)\) be the registration record of a transition \(u:g\to g'\).
The transition is Grand-safe when every changed port is registered:
\[
  \{q\in F_6\mid g(q)\ne g'(q)\}
  \subseteq
  \Reg(u).
\]
It is Grand-unsafe when a changed port is absent from \(\Reg(u)\).
\end{definition}

Equivalently,
\[
  \text{unsafe}=\text{unregistered discontinuity}.
\]
The discontinuity itself is not the violation.  The violation is that the
boundary crossing occurred while the theory failed to enter it into the
ledger.

\section{Conclusion}

The Grand Meson is not proposed as an omnipotent meson.  It is proposed as
a safe \(n=3\) general-purpose meson:
\[
  \mathsf{Grand}=\mathsf{Pol}^{6}.
\]
It conservatively recovers the General Meson by projection, replaces
undefined ports with typed ports, avoids an absolute aether frame, and
treats safety as registration rather than immobility.

The path is therefore narrow.  Begin with the \(\qu\leftrightarrow\qd\)
eye that already works.  Open \(\qc\leftrightarrow\qs\) and
\(\qt\leftrightarrow\qb\) without breaking it.  Stop at \(n=3\) until a
future theory can make higher axes readable without turning them into a
burden of infinite explanation.

\begin{thebibliography}{9}

\bibitem{Yukawa1935}
H. Yukawa.
On the interaction of elementary particles.
\emph{Proceedings of the Physico-Mathematical Society of Japan},
17:48--57, 1935.

\bibitem{Einstein1905}
A. Einstein.
On the electrodynamics of moving bodies.
\emph{Annalen der Physik}, 17:891--921, 1905.

\end{thebibliography}

\end{document}
